\section{Completion and the geometry of boundary strata}\label{sec:geometry}

\subsection{Completion and finite curvature energy}

\begin{theorem}\label{thm:completion}
The metric completion of $(\M,\gL)$ is canonically homeomorphic to $\MU$.
It is the compact length space defined by the canonical metric on $\Mn$.
It is a noncollapsed $\operatorname{RCD}(0,4N)$ space, with its Riemannian
volume measure. If $\eta$ is a K\"ahler class on $\MG$, the Ricci-flat
metrics in $\Lambda+\varepsilon\eta$ converge to this completion in the
Gromov--Hausdorff topology as $\varepsilon\downarrow0$, and converge
smoothly on compact subsets over $\Mn_{\rm reg}$.
\end{theorem}
\begin{proof}
The class $\Lambda'$ is a positive multiple of an ample rational class and
$\pi$ is a crepant resolution. Song's theorem
\cite[Theorem~1.1 and Corollary~1.1]{Song14} identifies the completion of
the regular metric with $\Mn$, proves compactness and geodesic convexity
of its regular set, and bounds the Hausdorff dimension of the singular
set by $4N-4$. The additional closed analytic set
$\Mn_{\rm reg}\setminus\M$ has real codimension at least four.
Removing it does not change intrinsic distances: a piecewise smooth curve
in the regular locus can be perturbed, with arbitrarily small length
change, to avoid an analytic subset of real codimension at least two.
It follows that completing $\M$ gives the same space.
Proposition~\ref{prop:algebraic-model} identifies its topology with $\MU$.

The convergence assertion is
\cite[Theorem~1.6]{CT15}, applied to the nef and big class $\Lambda$;
see also \cite[Theorems~1.1--1.2]{RZ11}.
The approximating metrics are Ricci-flat and their volumes tend to the
positive number in \eqref{eq:volume}. Their compact Gromov--Hausdorff limit
is therefore noncollapsed. Stability of the synthetic Ricci lower bound
and volume convergence give the stated RCD structure
\cite{DPG18}. Here noncollapsed means that the reference measure is the
$4N$-dimensional Hausdorff measure, with its usual Euclidean normalization.
\end{proof}

\begin{corollary}\label{cor:bounds}
Writing $V=\Vol(X,g_X)$ and $D=\operatorname{diam}(\MU)$, one has
\begin{equation}\label{eq:diameter}
 D\ge (16\pi^2V)^{1/4}
       \left(\frac{(2N)!}{N!}\right)^{1/(4N)}
   \ge \left(\frac{64\pi^2NV}{e}\right)^{1/4}.
\end{equation}
The curvature is square integrable on $\M$. More precisely, define
$|R|_{\rm op}^2=\frac14\sum_{i,j,k,l}R_{ijkl}^2$ in a real orthonormal
frame. Then
\begin{equation}\label{eq:curvature-bound}
 \int_{\M}|R|_{\rm op}^2dV_{\gL}
 \le \frac{48\pi^2(N+3)}{(N-1)!}(16\pi^4V)^{N-1}.
\end{equation}
\end{corollary}
\begin{proof}
Bishop--Gromov comparison gives
$\Vol(\MU)\le\omega_{4N}D^{4N}$, where
$\omega_{4N}=\pi^{2N}/(2N)!$. Substitution of \eqref{eq:volume} gives
the first bound. Integrating $\log x$ over $[N,2N]$ yields
$(2N)!/N!\ge(4N/e)^N$, giving the second.

For a compact Ricci-flat K\"ahler manifold of complex dimension $2N$, the
Chern--Weil identity in this norm is
\[
 \int |R|_{\rm op}^2dV
 =\frac{8\pi^2}{(2N-2)!}\int c_2\,\Omega^{2N-2}
\]
\cite[formula following (8)]{HS01}.
Apply it to the approximating metrics in Theorem~\ref{thm:completion}
and use smooth convergence on compact subsets and exhaustion. The
right side tends to its value at $\Lambda$.
For K3$^{[N]}$ type,
\[
 \int c_2\alpha^{2N-2}
 =6(N+3)(2N-3)!!\,q(\alpha)^{N-1}.
\]
This coefficient follows as well by comparing the degree $N-1$ terms in
Riemann--Roch and
$\chi(L)=\binom{q(c_1(L))/2+N+1}{N}$.
Using $q(\Lambda)=32\pi^4V$ gives \eqref{eq:curvature-bound}.
\end{proof}

\subsection{Restriction of the current to a stratum}
Put $v_\ell=v+\ell(0,0,1)$ and let $\M_H(v_\ell)$ denote the stable
locally free locus, when nonempty. On the symmetric product, let
$\omega_X^{(\ell)}$ be the current whose pullback to $X^\ell$ is
$\sum_{a=1}^\ell\operatorname{pr}_a^*\omega$.

\begin{theorem}\label{thm:stratum}
There is a holomorphic map
\[
 j_\ell:\M_H(v_\ell)\times\Sym^\ell X\longrightarrow\Mn,
 \qquad (F,Z)\longmapsto(F,Z),
\]
which is injective on points. Its pullback current is
\begin{equation}\label{eq:stratum-current}
 j_\ell^*T=\Omega_{L^2,v_\ell}+4\pi^2\omega_X^{(\ell)}.
\end{equation}
This identity includes cycles with repeated points.
\end{theorem}
\begin{proof}
Work over a local universal family of stable bundles $F$.
The relative Quot scheme of length-$\ell$ quotients of $F$ maps to $\MG$
by taking kernels and to $\Sym^\ell X$ by their cycles. The kernels are
slope stable. They belong to the unique irreducible Gieseker moduli space
$\MG$, hence to the closure of $\M$. Fibres over a cycle are products of
punctual Quot schemes and are connected by
\cite[Theorem~1]{EL99}. Composition with $\pi$ is constant on each such
fibre. The base $\M_H(v_\ell)\times\Sym^\ell X$ is normal; proper descent
therefore gives the required holomorphic map. These local constructions
agree, including when the universal families differ by a line from the
base. The description of points gives injectivity.

Fix $(F,Z)$ and choose $C,G$ as in Section~\ref{sec:continuity}.
In \eqref{eq:ideal-potential}, the first two terms depend only on $F$ and
have Hessian $\Omega_{L^2,v_\ell}$. On a neighbourhood disjoint from $C$ the
last term has Hessian
\[
 \ddc\left(-4\pi^2\sum_an_ag_C(x_a)\right)
       =4\pi^2\sum_an_a\omega(x_a).
\]
This proves \eqref{eq:stratum-current} on the ordered product; the
symmetric expression descends, also across its diagonals. All potentials
are continuous, so their pullbacks and Hessians are defined there.
\end{proof}

The current identity alone does not identify an induced Riemannian metric
on a singular stratum. The next theorem supplies that extra regularity in
rank two with distinct bubble points.

\subsection{Rank two and distinct points}
Let $r=2$, and let $S_\ell$ be the image of the locus with stable $F$ and
$\ell$ distinct points, each of multiplicity one.

\begin{theorem}\label{thm:orbifold}
Near $S_\ell$, the completed metric is a smooth hyperk\"ahler orbifold
metric. The stratum is totally geodesic, and its intrinsic metric is
\begin{equation}\label{eq:stratum-metric}
 \gL{}_{,v_\ell}+4\pi^2g_X^{(\ell)}.
\end{equation}
At $p\in S_\ell$ the metric tangent cone and normalized volume density are
\begin{equation}\label{eq:orbifold-cone}
 \RR^{4N-4\ell}\times(\RR^4/\{\pm1\})^\ell,
 \qquad \vartheta(p)=2^{-\ell}.
\end{equation}
In equivariant normal coordinates on the covering chart, the metric
coefficients differ from their values at the centre by $O(\rho^2)$.
\end{theorem}
\begin{proof}
The ADHM space of one rank-two instanton is
$\frM_0(2,1)=\CC^2\times(\CC^2/\{\pm1\})$.
Equation~\eqref{eq:adhm-germ} gives the local quotient
\[
 \CC^{2N-2\ell}\times(\CC^2/\{\pm1\})^\ell.
\]
The orbifold regularity theorem of Guenancia--Pan--P\u aun
\cite[Corollary~B]{GPP25} says that the canonical Ricci-flat metric is
smooth on the whole quotient-singularity locus. Its pullback to the
covering chart is thus a smooth Ricci-flat metric. The parallel symplectic
form and complex structures extend over the fixed loci, so the metric is
hyperk\"ahler there as well.

The stratum lifts to the common fixed-point locus of the isometric group
$(\ZZ/2)^\ell$. A fixed-point locus of isometries is totally geodesic.
Its metric is \eqref{eq:stratum-metric}, since its K\"ahler form is the
restriction computed in Theorem~\ref{thm:stratum}.
The invariant tangent-space decomposition gives \eqref{eq:orbifold-cone};
the finite quotient divides Euclidean volume by $2^\ell$.
Finally the first derivatives of a smooth metric vanish in normal
coordinates. The exponential map is equivariant, giving the stated
quadratic coefficient estimate.
\end{proof}

In particular, if the lower moduli space in $S_1$ satisfies the hypotheses
of Theorem~\ref{thm:canonical}, then
\begin{equation}\label{eq:s1-volume}
 \Vol(S_1)=\frac{(16\pi^4V)^{N-1}}{(N-2)!}.
\end{equation}
Indeed its complex dimension is $2N-4$, and the product in
\eqref{eq:stratum-metric} has volume
\[\Vol(\M_H(v_1))\Vol(X,4\pi^2g_X).\]
The same calculation applies if the lower space consists of a single
spherical bundle, with volume one.
Equation~\eqref{eq:s1-volume} concerns $S_1$. It equals the
$(4N-4)$-dimensional measure of the entire singular locus only if its
complement there has measure zero in that dimension. Reducible residual
bundles require a separate analysis; odd first Chern class alone does not
supply that conclusion.

\subsection{A Lipschitz estimate for curvature measures}
The following estimate does not require a complex structure on the base.

\begin{proposition}\label{prop:lipschitz}
Let $X$ be a compact oriented Riemannian four-manifold and let $\M$ be a
smooth moduli space of irreducible projectively anti-self-dual connections
with fixed determinant and energy $8\pi^2\kappa$. For every Lipschitz
function $f$ on $X$, the function
\[
 \Phi_f(A)=\frac12\int_X f|F_A^0|^2dV_X
\]
is $2\pi\sqrt\kappa\,\operatorname{Lip}(f)$-Lipschitz in the $L^2$ metric.
In the K3 setting it extends to the completion by
\begin{equation}\label{eq:phi-ideal}
 \Phi_f(F,Z)=\frac12\int_Xf|F_{A_F}^0|^2dV_X
                   +4\pi^2\sum_jn_jf(x_j).
\end{equation}
In particular, two one-bubble points with the same background obey
\begin{equation}\label{eq:bubble-distance}
 d_{\MU}((F,x),(F,y))\ge\frac{2\pi}{\sqrt\kappa}d_X(x,y).
\end{equation}
\end{proposition}
\begin{proof}
For smooth $f$ and a harmonic tangent vector $a$, the Yang--Mills equation
and integration by parts give
\[
 d\Phi_f(a)=-\int_X\langle\iota_{\nabla f}F_A^0,a\rangle dV_X.
\]
An anti-self-dual two-form satisfies
$|\iota_wF_A^0|^2=\tfrac12|w|^2|F_A^0|^2$.
Cauchy--Schwarz therefore bounds the absolute value by
$2\pi\sqrt\kappa\|df\|_\infty\|a\|_{L^2}$.
The same integration by parts holds for a Lipschitz $f$ using its weak
first derivatives, or follows by local smooth approximation. Integration
along paths proves the asserted Lipschitz constant.
The completion theorem and weak convergence of the energy measures give
\eqref{eq:phi-ideal}. For equal backgrounds its difference is
$4\pi^2(f(x)-f(y))$. Taking the supremum over all $1$-Lipschitz functions
proves \eqref{eq:bubble-distance}.
\end{proof}

For the rank-two stratum with a fixed stable background,
Theorem~\ref{thm:orbifold} also gives the upper bound
$d_{\MU}((F,x),(F,y))\le2\pi d_X(x,y)$ by paths inside the stratum.
For sufficiently close points in a convex orbifold neighbourhood this
upper bound is an equality. Global total geodesicity does not by itself
make every long intrinsic geodesic globally minimizing.
